% ============================================================
\chapter{Entra\^inement \`a Grande \'Echelle --- Distributed Training}
\label{ch:entrainement}
\index{entra\^inement distribu\'e}
% ============================================================

\section{Pourquoi l'entra\^inement distribu\'e~?}
\label{sec:why-distributed}

Les mod\`eles modernes de deep learning atteignent des tailles qui rendent
l'entra\^inement sur un seul GPU impossible ou prohibitivement lent~:
\begin{itemize}
  \item \textbf{GPT-3}~: 175 milliards de param\`etres
  \item \textbf{LLaMA-2 70B}~: 70 milliards de param\`etres
  \item \textbf{Vision Transformers}~: entra\^inement sur des millions
    d'images en haute r\'esolution
\end{itemize}

L'entra\^inement distribu\'e r\'epartit le travail sur plusieurs GPU
(ou plusieurs machines) pour r\'eduire le temps d'entra\^inement et
permettre l'entra\^inement de mod\`eles qui ne tiennent pas en m\'emoire
sur un seul acc\'el\'erateur.

\subsection{Parall\'elisme de donn\'ees vs parall\'elisme de mod\`ele}
\index{parall\'elisme de donn\'ees}\index{parall\'elisme de mod\`ele}

\begin{definition}[Parall\'elisme de donn\'ees]
Chaque GPU poss\`ede une \textbf{copie compl\`ete} du mod\`ele. Le batch
de donn\'ees est divis\'e en sous-batches, chaque GPU traite un sous-batch,
puis les gradients sont agr\'eg\'es (moyenne) avant la mise \`a jour des
poids.
\end{definition}

\begin{definition}[Parall\'elisme de mod\`ele]
Le mod\`ele est \textbf{partitionn\'e} entre plusieurs GPU. Chaque GPU
ne contient qu'une partie du mod\`ele (par couches ou par tenseurs). Les
activations sont transmises entre GPU lors du forward pass.
\end{definition}

\begin{center}
\begin{tikzpicture}[
  gpu/.style={rectangle, draw, rounded corners, minimum width=2.2cm,
    minimum height=1.8cm, align=center, font=\scriptsize},
  batch/.style={rectangle, draw, fill=blue!10, minimum width=1.8cm,
    minimum height=0.5cm, font=\scriptsize, rounded corners},
  model/.style={rectangle, draw, fill=green!10, minimum width=1.8cm,
    minimum height=0.5cm, font=\scriptsize, rounded corners},
  arr/.style={-{Stealth[length=2mm]}, thick},
  title/.style={font=\small\bfseries}
]
  % Data Parallelism
  \node[title] at (0, 3) {Parall\'elisme de donn\'ees};

  \node[gpu, fill=yellow!5] (g1) at (-1.5, 1) {};
  \node[model] at ([yshift=0.3cm]g1.center) {Mod\`ele complet};
  \node[batch, fill=red!15] at ([yshift=-0.3cm]g1.center) {Batch 1/3};
  \node[below=0.05cm, font=\tiny] at (g1.south) {GPU 0};

  \node[gpu, fill=yellow!5] (g2) at (1.5, 1) {};
  \node[model] at ([yshift=0.3cm]g2.center) {Mod\`ele complet};
  \node[batch, fill=red!15] at ([yshift=-0.3cm]g2.center) {Batch 2/3};
  \node[below=0.05cm, font=\tiny] at (g2.south) {GPU 1};

  \node[gpu, fill=yellow!5] (g3) at (4.5, 1) {};
  \node[model] at ([yshift=0.3cm]g3.center) {Mod\`ele complet};
  \node[batch, fill=red!15] at ([yshift=-0.3cm]g3.center) {Batch 3/3};
  \node[below=0.05cm, font=\tiny] at (g3.south) {GPU 2};

  \node[rectangle, draw, dashed, fill=orange!5, minimum width=3cm,
    minimum height=0.5cm, font=\scriptsize, rounded corners,
    below=0.6cm of g2] (sync) {AllReduce (gradients)};
  \draw[arr, dashed, gray] (g1.south) -- (sync.west);
  \draw[arr, dashed, gray] (g2.south) -- (sync.north);
  \draw[arr, dashed, gray] (g3.south) -- (sync.east);

  % Model Parallelism
  \node[title] at (10, 3) {Parall\'elisme de mod\`ele};

  \node[gpu, fill=yellow!5] (m1) at (8.5, 1) {};
  \node[model, fill=green!20] at ([yshift=0.3cm]m1.center) {Couches 1--4};
  \node[batch] at ([yshift=-0.3cm]m1.center) {Batch complet};
  \node[below=0.05cm, font=\tiny] at (m1.south) {GPU 0};

  \node[gpu, fill=yellow!5] (m2) at (11.5, 1) {};
  \node[model, fill=green!30] at ([yshift=0.3cm]m2.center) {Couches 5--8};
  \node[batch] at ([yshift=-0.3cm]m2.center) {Activations};
  \node[below=0.05cm, font=\tiny] at (m2.south) {GPU 1};

  \draw[arr, steelblue] (m1.east) -- (m2.west)
    node[midway, above, font=\tiny] {forward};
  \draw[arr, deepred, bend right=20] (m2.south west) -- (m1.south east)
    node[midway, below, font=\tiny] {backward};
\end{tikzpicture}
\end{center}

\section{Formulation math\'ematique du SGD parall\`ele}
\label{sec:math-parallel-sgd}
\index{SGD parall\`ele}

Soit $N$ le nombre de GPU (workers). \`A chaque it\'eration~:

\begin{formules}[title=\textbf{SGD avec parall\'elisme de donn\'ees}]
\begin{enumerate}
  \item Le batch global $\mathcal{B}$ de taille $B$ est divis\'e en
    $N$ sous-batches $\mathcal{B}_k$ de taille $B/N$.
  \item Chaque worker $k$ calcule son gradient local~:
    \[
      g_k = \frac{1}{|\mathcal{B}_k|} \sum_{(x,y) \in \mathcal{B}_k}
      \nabla_\theta \, \ell(f_\theta(x), y)
    \]
  \item Les gradients sont agr\'eg\'es par AllReduce~:
    \[
      \bar{g} = \frac{1}{N} \sum_{k=1}^{N} g_k
    \]
  \item Chaque worker met \`a jour ses poids~:
    \[
      \theta \leftarrow \theta - \eta \, \bar{g}
    \]
\end{enumerate}
Le gradient agr\'eg\'e $\bar{g}$ est un estimateur non biais\'e du gradient
sur le batch complet. Le batch effectif est de taille $B = N \cdot B_{\text{local}}$.
\end{formules}

\begin{attention}[title=\textbf{Scaling du learning rate}]
Avec $N$ GPU, le batch effectif est multipli\'e par $N$. La r\`egle
empirique (\emph{linear scaling rule})~:
\[
  \eta_{\text{distribu\'e}} = N \cdot \eta_{\text{base}}
\]
En pratique, on utilise un \textbf{warmup lin\'eaire} sur les premi\`eres
it\'erations pour stabiliser l'entra\^inement avec un grand learning rate.
\end{attention}

\subsection{AllReduce}
\index{AllReduce}

L'op\'eration \textbf{AllReduce} est le c\oe ur de la synchronisation
des gradients. Elle agr\`ege les tenseurs de tous les workers et distribue
le r\'esultat \`a chacun.

\begin{center}
\begin{tikzpicture}[
  gpu/.style={rectangle, draw, rounded corners, minimum width=1.4cm,
    minimum height=0.8cm, align=center, font=\scriptsize, fill=yellow!10},
  arr/.style={-{Stealth[length=2mm]}, thick, steelblue}
]
  % Before
  \node[font=\small\bfseries] at (0, 2.2) {Avant AllReduce};
  \node[gpu] (b0) at (-1.5, 1) {$g_0$};
  \node[gpu] (b1) at (0, 1) {$g_1$};
  \node[gpu] (b2) at (1.5, 1) {$g_2$};

  % After
  \node[font=\small\bfseries] at (5.5, 2.2) {Apr\`es AllReduce};
  \node[gpu] (a0) at (4, 1) {$\bar{g}$};
  \node[gpu] (a1) at (5.5, 1) {$\bar{g}$};
  \node[gpu] (a2) at (7, 1) {$\bar{g}$};

  \draw[arr] (2.2, 1) -- (3.3, 1)
    node[midway, above, font=\scriptsize] {AllReduce};

  % Ring
  \node[font=\small\bfseries] at (3.5, -0.5) {Ring AllReduce};
  \node[gpu] (r0) at (1.5, -2) {GPU 0};
  \node[gpu] (r1) at (3.5, -1) {GPU 1};
  \node[gpu] (r2) at (5.5, -2) {GPU 2};
  \node[gpu] (r3) at (3.5, -3) {GPU 3};

  \draw[arr] (r0) -- (r1);
  \draw[arr] (r1) -- (r2);
  \draw[arr] (r2) -- (r3);
  \draw[arr] (r3) -- (r0);
\end{tikzpicture}
\end{center}

\section{PyTorch DistributedDataParallel (DDP)}
\label{sec:pytorch-ddp}
\index{DistributedDataParallel}\index{DDP}

\textbf{DDP}\index{DDP} est le m\'ecanisme standard de PyTorch pour le
parall\'elisme de donn\'ees. Chaque processus g\`ere un GPU, et les
gradients sont synchronis\'es automatiquement via NCCL\index{NCCL}.

\begin{codeblock}[title=\textbf{Entra\^inement DDP complet}]
\begin{minted}{python}
"""Distributed Data Parallel training with PyTorch."""
import os
import torch
import torch.nn as nn
import torch.optim as optim
import torch.distributed as dist
from torch.nn.parallel import DistributedDataParallel as DDP
from torch.utils.data import DataLoader, DistributedSampler
from torchvision import datasets, transforms


def setup(rank: int, world_size: int) -> None:
    """Initialize the distributed process group."""
    os.environ["MASTER_ADDR"] = "localhost"
    os.environ["MASTER_PORT"] = "12355"
    dist.init_process_group("nccl", rank=rank, world_size=world_size)
    torch.cuda.set_device(rank)


def cleanup() -> None:
    """Destroy the process group."""
    dist.destroy_process_group()


class SimpleCNN(nn.Module):
    def __init__(self, num_classes: int = 10):
        super().__init__()
        self.features = nn.Sequential(
            nn.Conv2d(1, 32, 3, padding=1), nn.ReLU(),
            nn.MaxPool2d(2),
            nn.Conv2d(32, 64, 3, padding=1), nn.ReLU(),
            nn.MaxPool2d(2),
        )
        self.classifier = nn.Sequential(
            nn.Flatten(),
            nn.Linear(64 * 7 * 7, 128), nn.ReLU(),
            nn.Linear(128, num_classes),
        )

    def forward(self, x):
        return self.classifier(self.features(x))


def train(rank: int, world_size: int, epochs: int = 10) -> None:
    """Training loop for one process (one GPU)."""
    setup(rank, world_size)

    # Data
    transform = transforms.Compose([
        transforms.ToTensor(),
        transforms.Normalize((0.1307,), (0.3081,)),
    ])
    dataset = datasets.MNIST(
        "data", train=True, download=True, transform=transform
    )
    sampler = DistributedSampler(
        dataset, num_replicas=world_size, rank=rank, shuffle=True
    )
    loader = DataLoader(
        dataset, batch_size=64, sampler=sampler, num_workers=2,
        pin_memory=True,
    )

    # Model
    model = SimpleCNN().to(rank)
    model = DDP(model, device_ids=[rank])

    # Optimizer
    optimizer = optim.Adam(model.parameters(), lr=1e-3)
    criterion = nn.CrossEntropyLoss()

    # Training loop
    for epoch in range(epochs):
        sampler.set_epoch(epoch)  # important for shuffling
        model.train()
        total_loss = 0.0

        for batch_idx, (data, target) in enumerate(loader):
            data, target = data.to(rank), target.to(rank)

            optimizer.zero_grad()
            output = model(data)
            loss = criterion(output, target)
            loss.backward()
            optimizer.step()

            total_loss += loss.item()

        if rank == 0:
            avg_loss = total_loss / len(loader)
            print(f"Epoch {epoch+1}/{epochs}, Loss: {avg_loss:.4f}")

    # Save model (only on rank 0)
    if rank == 0:
        torch.save(model.module.state_dict(), "model_ddp.pt")

    cleanup()


if __name__ == "__main__":
    world_size = torch.cuda.device_count()
    torch.multiprocessing.spawn(
        train, args=(world_size,), nprocs=world_size, join=True
    )
\end{minted}
\end{codeblock}

\begin{codeblock}[title=\textbf{Lancement avec torchrun}]
\begin{minted}{bash}
# Sur une seule machine avec 4 GPU
torchrun --nproc_per_node=4 train_ddp.py

# Sur 2 machines avec 4 GPU chacune
# Machine 0 (master)
torchrun --nproc_per_node=4 --nnodes=2 --node_rank=0 \
    --master_addr=192.168.1.1 --master_port=12355 train_ddp.py

# Machine 1
torchrun --nproc_per_node=4 --nnodes=2 --node_rank=1 \
    --master_addr=192.168.1.1 --master_port=12355 train_ddp.py
\end{minted}
\end{codeblock}

\begin{bonnepratique}[title=\textbf{Bonnes pratiques DDP}]
\begin{enumerate}
  \item \textbf{Appeler \cli{sampler.set\_epoch(epoch)}}~: sinon le
    shuffling sera identique \`a chaque \'epoque.
  \item \textbf{Sauvegarder via \cli{model.module}}~: DDP wrappe le
    mod\`ele, il faut acc\'eder au module interne.
  \item \textbf{Sauvegarder uniquement sur rank 0}~: \'evite les
    \'ecritures concurrentes.
  \item \textbf{Utiliser \cli{pin\_memory=True}}~: acc\'el\`ere les
    transferts CPU $\to$ GPU.
  \item \textbf{Pr\'ef\'erer \cli{torchrun}}~: g\`ere automatiquement les
    variables d'environnement (\cli{RANK}, \cli{WORLD\_SIZE}).
\end{enumerate}
\end{bonnepratique}

\section{Multi-GPU avec Accelerate (Hugging Face)}
\label{sec:accelerate}
\index{Accelerate}

\textbf{Accelerate}\index{Accelerate} de Hugging Face abstrait la
complexit\'e de DDP derri\`ere une API minimale. Le m\^eme code fonctionne
sur CPU, un GPU, multi-GPU et multi-machines.

\begin{codeblock}[title=\textbf{Entra\^inement avec Accelerate}]
\begin{minted}{python}
"""Training with Hugging Face Accelerate."""
from accelerate import Accelerator
import torch
import torch.nn as nn
from torch.utils.data import DataLoader
from torchvision import datasets, transforms

accelerator = Accelerator(mixed_precision="fp16")

# Data
transform = transforms.Compose([
    transforms.ToTensor(),
    transforms.Normalize((0.1307,), (0.3081,)),
])
dataset = datasets.MNIST("data", train=True, download=True,
                          transform=transform)
loader = DataLoader(dataset, batch_size=64, shuffle=True)

# Model and optimizer
model = SimpleCNN()  # meme modele que precedemment
optimizer = torch.optim.Adam(model.parameters(), lr=1e-3)
criterion = nn.CrossEntropyLoss()

# Accelerate prepare: wraps model, optimizer, dataloader
model, optimizer, loader = accelerator.prepare(
    model, optimizer, loader
)

# Training loop (identique au code single-GPU !)
for epoch in range(10):
    model.train()
    total_loss = 0.0
    for data, target in loader:
        optimizer.zero_grad()
        output = model(data)
        loss = criterion(output, target)
        accelerator.backward(loss)  # seul changement
        optimizer.step()
        total_loss += loss.item()

    if accelerator.is_main_process:
        print(f"Epoch {epoch+1}, Loss: {total_loss/len(loader):.4f}")

# Save
accelerator.wait_for_everyone()
if accelerator.is_main_process:
    unwrapped = accelerator.unwrap_model(model)
    torch.save(unwrapped.state_dict(), "model_accelerate.pt")
\end{minted}
\end{codeblock}

\begin{codeblock}[title=\textbf{Configuration et lancement Accelerate}]
\begin{minted}{bash}
# Installer
pip install accelerate

# Configuration interactive
accelerate config

# Lancer l'entrainement
accelerate launch --num_processes=4 train_accelerate.py

# Ou avec un fichier de config
accelerate launch --config_file accelerate_config.yaml \
    train_accelerate.py
\end{minted}
\end{codeblock}

\section{Entra\^inement en pr\'ecision mixte (AMP)}
\label{sec:mixed-precision}
\index{pr\'ecision mixte}\index{AMP}

L'entra\^inement en \textbf{pr\'ecision mixte}\index{pr\'ecision mixte}
utilise FP16 (ou BF16) pour les calculs rapides et FP32 pour les
op\'erations sensibles (accumulation de gradients, normalisation). Cela
r\'eduit la m\'emoire GPU de $\sim$50\% et acc\'el\`ere l'entra\^inement
de 1.5--3$\times$ sur les GPU modernes (Ampere, Hopper).

\begin{codeblock}[title=\textbf{AMP avec PyTorch natif}]
\begin{minted}{python}
"""Mixed precision training with torch.amp."""
from torch.amp import autocast, GradScaler

scaler = GradScaler("cuda")

for epoch in range(epochs):
    for data, target in loader:
        data, target = data.to(device), target.to(device)

        optimizer.zero_grad()

        # Forward pass in FP16
        with autocast("cuda"):
            output = model(data)
            loss = criterion(output, target)

        # Backward pass with gradient scaling
        scaler.scale(loss).backward()
        scaler.step(optimizer)
        scaler.update()
\end{minted}
\end{codeblock}

\begin{remarque}
Le \cli{GradScaler} est n\'ecessaire pour \'eviter les \emph{underflows}
en FP16~: les gradients tr\`es petits deviennent z\'ero en demi-pr\'ecision.
Le scaler multiplie la loss par un facteur \'elev\'e, puis divise les
gradients avant la mise \`a jour. Avec BF16 (disponible sur Ampere+), le
scaler n'est g\'en\'eralement pas n\'ecessaire car la plage dynamique de
BF16 est \'equivalente \`a celle de FP32.
\end{remarque}

\section{Optimisation des hyperparam\`etres avec Optuna}
\label{sec:optuna}
\index{Optuna}\index{hyperparam\`etres}

\textbf{Optuna}\index{Optuna} est une biblioth\`eque d'optimisation
d'hyperparam\`etres qui utilise des algorithmes bayesiens (TPE) pour
explorer efficacement l'espace de recherche.

\begin{codeblock}[title=\textbf{Optimisation avec Optuna}]
\begin{minted}{python}
"""Hyperparameter optimization with Optuna."""
import optuna
import torch
import torch.nn as nn
from torch.utils.data import DataLoader
from sklearn.metrics import accuracy_score


def objective(trial: optuna.Trial) -> float:
    """Objective function for Optuna."""
    # Suggest hyperparameters
    lr = trial.suggest_float("lr", 1e-5, 1e-2, log=True)
    n_layers = trial.suggest_int("n_layers", 1, 4)
    hidden_size = trial.suggest_categorical(
        "hidden_size", [64, 128, 256, 512]
    )
    dropout = trial.suggest_float("dropout", 0.0, 0.5)
    batch_size = trial.suggest_categorical(
        "batch_size", [32, 64, 128]
    )

    # Build model dynamically
    layers = []
    in_features = 784  # MNIST
    for i in range(n_layers):
        layers.extend([
            nn.Linear(in_features, hidden_size),
            nn.ReLU(),
            nn.Dropout(dropout),
        ])
        in_features = hidden_size
    layers.append(nn.Linear(in_features, 10))
    model = nn.Sequential(*layers).to("cuda")

    optimizer = torch.optim.Adam(model.parameters(), lr=lr)
    criterion = nn.CrossEntropyLoss()

    # Train for a few epochs
    for epoch in range(5):
        model.train()
        for data, target in train_loader:
            data = data.view(data.size(0), -1).to("cuda")
            target = target.to("cuda")
            optimizer.zero_grad()
            loss = criterion(model(data), target)
            loss.backward()
            optimizer.step()

        # Pruning: report intermediate value
        val_acc = evaluate(model, val_loader)
        trial.report(val_acc, epoch)
        if trial.should_prune():
            raise optuna.TrialPruned()

    return val_acc


# Run optimization
study = optuna.create_study(
    direction="maximize",
    pruner=optuna.pruners.MedianPruner(n_warmup_steps=2),
)
study.optimize(objective, n_trials=50, timeout=3600)

# Results
print(f"Best trial: {study.best_trial.value:.4f}")
print(f"Best params: {study.best_trial.params}")

# Visualization
optuna.visualization.plot_optimization_history(study)
optuna.visualization.plot_param_importances(study)
\end{minted}
\end{codeblock}

\begin{bonnepratique}[title=\textbf{Bonnes pratiques Optuna}]
\begin{enumerate}
  \item \textbf{Utiliser le pruning}~: arr\^eter les essais peu prometteurs
    t\^ot pour \'economiser du temps de calcul.
  \item \textbf{Espaces log-uniformes pour les learning rates}~: utiliser
    \cli{suggest\_float(..., log=True)} pour les param\`etres \`a \'echelle
    logarithmique.
  \item \textbf{Stocker les r\'esultats}~: utiliser un backend persistant
    (SQLite, MySQL) pour reprendre une \'etude interrompue.
  \item \textbf{Parall\'eliser les essais}~: lancer plusieurs workers
    Optuna en parall\`ele avec une base de donn\'ees partag\'ee.
\end{enumerate}
\end{bonnepratique}

\section{Comparaison des frameworks d'entra\^inement distribu\'e}
\label{sec:distributed-comparison}

\begin{center}
\begin{tabular}{@{}lcccc@{}}
\toprule
\textbf{Crit\`ere} & \textbf{DDP} & \textbf{FSDP} &
  \textbf{DeepSpeed} & \textbf{Horovod} \\
\midrule
Type & Data parallel & Data + model & Data + model & Data parallel \\
Framework & PyTorch & PyTorch & PyTorch/TF & PyTorch/TF \\
Sharding poids & Non & Oui & Oui (ZeRO) & Non \\
Offload CPU & Non & Oui & Oui & Non \\
Facilit\'e & Moyen & Moyen & Difficile & Facile \\
Cas d'usage & Mod\`eles $<$10B & 10--100B & 10--1000B & Clusters HPC \\
Optimiseur int\'egr\'e & Non & Non & Oui & Non \\
Mixed precision & Manuel & Int\'egr\'e & Int\'egr\'e & Manuel \\
\bottomrule
\end{tabular}
\end{center}

\begin{attention}[title=\textbf{Choisir le bon framework}]
\begin{itemize}
  \item \textbf{DDP}~: d\'efaut pour tout mod\`ele qui tient sur un GPU.
    Simple et performant.
  \item \textbf{FSDP}~: mod\`eles trop grands pour un GPU mais qui tiennent
    en m\'emoire agr\'eg\'ee. Solution native PyTorch.
  \item \textbf{DeepSpeed}~: mod\`eles tr\`es grands (LLMs). ZeRO Stage 3
    permet d'entra\^iner des mod\`eles de centaines de milliards de
    param\`etres. Configuration plus complexe.
  \item \textbf{Horovod}~: populaire dans les clusters HPC avec MPI.
    API simple (\cli{hvd.init()}, \cli{hvd.DistributedOptimizer()}).
\end{itemize}
\end{attention}

\begin{center}
\begin{tikzpicture}[
  zone/.style={rectangle, draw, rounded corners, minimum height=1.2cm,
    align=center, font=\small},
  arr/.style={-{Stealth[length=2mm]}, thick, steelblue}
]
  \node[zone, fill=green!10, minimum width=3cm] (small) at (0, 0)
    {\textbf{$<$1B params}\\DDP};
  \node[zone, fill=yellow!10, minimum width=3cm] (medium) at (4, 0)
    {\textbf{1--10B params}\\FSDP / DeepSpeed\\ZeRO Stage 2};
  \node[zone, fill=orange!10, minimum width=3cm] (large) at (8.5, 0)
    {\textbf{10--100B params}\\DeepSpeed ZeRO 3\\+ offload CPU};
  \node[zone, fill=red!10, minimum width=3cm] (xlarge) at (13, 0)
    {\textbf{$>$100B params}\\Pipeline + Tensor\\parallelism};

  \draw[arr] (small) -- (medium);
  \draw[arr] (medium) -- (large);
  \draw[arr] (large) -- (xlarge);

  \node[font=\scriptsize, below=0.3cm of small] {1--8 GPU};
  \node[font=\scriptsize, below=0.3cm of medium] {8--32 GPU};
  \node[font=\scriptsize, below=0.3cm of large] {32--256 GPU};
  \node[font=\scriptsize, below=0.3cm of xlarge] {256+ GPU};
\end{tikzpicture}
\end{center}

\section{Mini-projet~: entra\^inement distribu\'e}
\label{sec:mini-project-ch6}

\begin{miniprojet}[title=\textbf{Mini-projet 6~: Entra\^inement multi-GPU}]
Mettez en place un pipeline d'entra\^inement distribu\'e~:
\begin{enumerate}
  \item \textbf{Baseline single-GPU}~: entra\^inez un mod\`ele CNN sur
    CIFAR-10 avec PyTorch standard. Mesurez le temps par \'epoque.
  \item \textbf{DDP}~: adaptez le code pour \cli{DistributedDataParallel}.
    Comparez le temps par \'epoque avec 2 et 4 GPU.
  \item \textbf{Mixed precision}~: ajoutez l'AMP (\cli{autocast} +
    \cli{GradScaler}). Mesurez le gain en m\'emoire et en vitesse.
  \item \textbf{Accelerate}~: r\'e\'ecrivez le code avec
    \pkg{accelerate}. V\'erifiez que les r\'esultats sont identiques.
  \item \textbf{Optuna}~: optimisez le learning rate, le batch size et
    l'architecture avec Optuna (20 trials minimum).
  \item \textbf{Rapport}~: pr\'esentez un tableau comparatif~: temps par
    \'epoque, m\'emoire GPU, accuracy finale pour chaque configuration.
\end{enumerate}
Livrable~: d\'ep\^ot Git avec les scripts, les r\'esultats et le rapport.
\end{miniprojet}

\section{Exercices}
\label{sec:exercises-ch6}

\begin{exercice}[$\bigstar$ --- Premier entra\^inement DDP]
Adaptez un script d'entra\^inement single-GPU existant pour utiliser DDP~:
\begin{enumerate}
  \item Ajoutez \cli{init\_process\_group} et \cli{DistributedSampler}
  \item Wrappez le mod\`ele avec \cli{DistributedDataParallel}
  \item Lancez avec \cli{torchrun --nproc\_per\_node=2}
  \item V\'erifiez que la loss converge de la m\^eme mani\`ere qu'en
    single-GPU
\end{enumerate}
\end{exercice}

\begin{exercice}[$\bigstar\bigstar$ --- Scaling du learning rate]
\'Etudiez l'impact du scaling du learning rate~:
\begin{enumerate}
  \item Entra\^inez un mod\`ele avec 1 GPU et $\eta = 0.001$, batch
    size $= 64$
  \item Passez \`a 4 GPU avec batch size effectif $= 256$~:
    \begin{itemize}
      \item Sans scaling du learning rate ($\eta = 0.001$)
      \item Avec scaling lin\'eaire ($\eta = 0.004$)
      \item Avec scaling lin\'eaire + warmup (5 \'epoques)
    \end{itemize}
  \item Tracez les courbes de loss et d'accuracy pour chaque configuration
  \item Concluez sur la n\'ecessit\'e du warmup
\end{enumerate}
\end{exercice}

\begin{exercice}[$\bigstar\bigstar\bigstar$ --- Pipeline complet avec Optuna et DDP]
Construisez un pipeline d'entra\^inement complet~:
\begin{enumerate}
  \item \'Ecrivez une fonction \cli{objective} Optuna qui entra\^ine un
    mod\`ele en DDP
  \item Optimisez au moins 5 hyperparam\`etres (learning rate, architecture,
    dropout, weight decay, scheduler)
  \item Utilisez le pruning m\'edian pour arr\^eter les essais peu
    prometteurs
  \item Stockez les r\'esultats dans une base SQLite pour la persistance
  \item Visualisez les r\'esultats avec les fonctions de visualisation
    Optuna (importance des param\`etres, historique, contour plots)
  \item R\'e-entra\^inez le meilleur mod\`ele avec AMP et sauvegardez-le
\end{enumerate}
\end{exercice}

\section{Aide-m\'emoire}

\begin{cheatsheet}[title=\textbf{R\'esum\'e du chapitre 6}]
\begin{tabular}{@{}p{3.5cm}p{9cm}@{}}
\toprule
\textbf{Concept} & \textbf{Point cl\'e} \\
\midrule
Data parallelism & Chaque GPU a le mod\`ele complet, batch divis\'e \\
Model parallelism & Mod\`ele partitionn\'e entre GPU \\
AllReduce & Agr\'egation des gradients~: $\bar{g} = \frac{1}{N}\sum g_k$ \\
DDP & \cli{DistributedDataParallel}, \cli{torchrun}, \cli{DistributedSampler} \\
Accelerate & \cli{accelerator.prepare()}, \cli{accelerator.backward()} \\
AMP & \cli{autocast("cuda")}, \cli{GradScaler}, $\sim$50\% m\'emoire \\
Linear scaling & $\eta_{\text{dist}} = N \cdot \eta_{\text{base}}$ + warmup \\
Optuna & \cli{study.optimize(objective, n\_trials)}, pruning, TPE \\
FSDP & Sharding des poids pour mod\`eles 10--100B \\
DeepSpeed & ZeRO Stages 1--3, offload CPU, mod\`eles $>$100B \\
\bottomrule
\end{tabular}
\end{cheatsheet}
